\chapter{Exercícios anexos}

\section{Exercícios de normalização}
Abaixo temos uma visão geral de informações mantidas, tradicionalmente em Excel, pela coordenação de estágios de um instituto universitário. Uma mesma empresa fornecedora de estágio pode ofertar vagas com orientadores distintos.

\textbf{Planilha de Estágio}:
(CNPJ\_Empresa, Nome\_Empresa, Endereco\_Empresa, Matrícula\_Aluno, Nome\_Aluno, Curso\_Aluno, Id\_Orientador, Nome\_Orientador, Titulacao\_Orientador, Data\_Inicio\_Estagio, Horas\_Mensais)

\textbf{Questão 1:} Descreva, com um exemplo em papel, quais tipos de Anomalia de Atualização e Anomalia de Exclusão poderiam ocorrer exatamente devido à estrutura acima, caso a empresa precisasse mudar de endereço em São Paulo.

\textbf{Questão 2:} Identifique os Grupos Repetitivos e liste qual seria a chave candidata composta ideal para normalizar o formulário diretamente para a sua Primeira Forma Normal (1FN). Apresente o resultado.

\textbf{Questão 3:} Conduza os dados gerados da 1FN aos moldes da Segunda Forma Normal (2FN). Quantas e quais tabelas resultaram deste processo de eliminação de dependências funcionais parciais? Elenque as chaves de tabela (PK e FK).

\textbf{Questão 4:} Prossiga com as tabelas geradas na Questão 3 e promova o projeto à sua Forma Normal de Boyce-Codd (FNBC), removendo quaisquer dependências transitivas. Especifique com clareza quais tabelas extras surgiram para alocar a entidade responsável. O seu Diagrama Final representará de fato uma modelagem íntegra e sem distorções operacionais.

\section{Exercícios de DDL (Data Definition Language)}

\subsection{Criação de Tabelas}
\begin{enumerate}
    \item Escreva o comando SQL para criar uma tabela \texttt{biblioteca} com as seguintes especificações:
    \begin{itemize}
        \item \texttt{id}: Inteiro, chave primária e autoincremento.
        \item \texttt{nome\_livro}: Texto (limite de 150 caracteres), preenchimento obrigatório.
        \item \texttt{ano\_lancamento}: Inteiro.
        \item \texttt{disponivel}: Booleano, com valor padrão verdadeiro.
    \end{itemize}

    \item Crie uma tabela chamada \texttt{autor} contendo:
    \begin{itemize}
        \item \texttt{id\_autor}: Inteiro, chave primária e autoincremento.
        \item \texttt{nome}: Texto (limite de 100 caracteres), preenchimento obrigatório.
        \item \texttt{nacionalidade}: Texto (limite de 50 caracteres), com valor padrão Brasileiro.
    \end{itemize}

    \item Crie uma tabela chamada \texttt{editora} contendo:
    \begin{itemize}
        \item \texttt{id\_editora}: Inteiro, chave primária e autoincremento.
        \item \texttt{nome\_fantasia}: Texto (limite de 100 caracteres), preenchimento obrigatório.
    \end{itemize}

    \item Crie uma tabela chamada \texttt{leitor} para registrar os usuários, contendo:
    \begin{itemize}
        \item \texttt{id\_leitor}: Inteiro, chave primária e autoincremento.
        \item \texttt{nome\_completo}: Texto (limite de 150 caracteres), preenchimento obrigatório.
        \item \texttt{data\_cadastro}: Data e hora, preenchimento obrigatório, com valor padrão a data e hora do sistema
    \end{itemize}

    \item Crie uma tabela de relacionamento chamada \texttt{emprestimo} contendo:
    \begin{itemize}
        \item \texttt{id\_emprestimo}: Inteiro, chave primária e autoincremento.
        \item \texttt{id\_livro}: Inteiro, chave estrangeira referenciando a tabela \texttt{biblioteca}.
        \item \texttt{id\_leitor}: Inteiro, chave estrangeira referenciando a tabela \texttt{leitor}.
        \item \texttt{data\_retirada}: Data e hora, preenchimento obrigatório, com valor padrão a data e hora do sistema.
    \end{itemize}
\end{enumerate}

\subsection{Alteração de Tabelas (\texttt{ALTER TABLE})}
\begin{enumerate}
    \item Adicione à tabela \texttt{biblioteca} a coluna \texttt{isbn}, definindo o tipo de dado para aceitar exatamente 13 caracteres.
    \item Modifique a coluna \texttt{nome\_livro} na tabela \texttt{biblioteca} para ampliar sua capacidade máxima para 250 caracteres.
    \item Adicione uma coluna \texttt{id\_autor} na tabela \texttt{biblioteca} e configure-a como chave estrangeira (\texttt{FOREIGN KEY}) referenciando a tabela \texttt{autor}.
    \item Adicione uma coluna \texttt{id\_editora} na tabela \texttt{biblioteca} e configure-a como chave estrangeira referenciando a tabela \texttt{editora}.
    \item Remova a coluna \texttt{ano\_lancamento} da tabela \texttt{biblioteca}.
\end{enumerate}

\subsection{Restrições e Exclusão}
\begin{enumerate}
    \item Adicione uma restrição (\texttt{CONSTRAINT}) de unicidade (\texttt{UNIQUE}) à coluna \texttt{isbn} da tabela \texttt{biblioteca}.
    \item Escreva o comando seguro para excluir a tabela \texttt{autor} garantindo que não ocorra erro caso ela não exista, mas não execute.
    \item Qual seria o impacto de tentar executar o comando do exercício anterior se já existirem livros cadastrados na tabela \texttt{biblioteca} associados a esses autores?
\end{enumerate}
